\chapter{Metabolic-Prior Replicator Dynamics for In-Vivo Oral Microbiome}
\label{ch:dieckow}

Where Chapter~\ref{ch:heine} identified interaction matrices from a defined
\emph{in-vitro} consortium, this chapter turns to the harder \emph{in-vivo}
regime: longitudinal 16S~rRNA data from a clinical cohort, where the community is
ten guilds wide, the time series only three points deep, and the interactions must
be regularised to be identifiable at all. The central idea of the thesis appears
here in full --- the off-diagonal \emph{sign} of the interaction matrix is
constrained from metabolic evidence (a flux-balance-derived sign prior), while its
magnitude is left to the data --- and is shown to be the property that actually
generalises out of sample.

\medskip\noindent\textit{This chapter is based on the internal analysis document:}
Nishioka K., Szafra\'{n}ski S.P., ``Metabolic-Prior Replicator Dynamics Predict
Oral Microbiome Community Composition under Leave-One-Out Cross-Validation''
(NIFE/IKM collaboration).

%==============================================================================
\section{Dataset and Guild Taxonomy}
\label{sec:dieckow_data}
%==============================================================================

We use the 16S~rRNA amplicon data of Dieckow et al.~\cite{dieckow2024}, a
prospective study of subgingival plaque on titanium implant abutments from
$10$ patients sampled at three weekly time points (W1, W2, W3). Reads were
processed by a standard amplicon pipeline --- ASV inference with DADA2~\cite{Callahan2016DADA2},
taxonomic classification against SILVA~\cite{Quast2013SILVA, Bokulich2018FeatureClassifier, Robeson2021RESCRIPt} within
QIIME2~\cite{Bolyen2019QIIME2} (VSEARCH~\cite{Rognes2016VSEARCH}), orchestrated by Nextflow~\cite{DiTommaso2017Nextflow}, referenced to the
human oral microbiome taxonomy~\cite{Dewhirst2010HOMD} --- and the resulting ASVs are
aggregated to the canonical $10$ class-level guilds used throughout this thesis
(Table~\ref{tab:dieckow_guilds}; the ordering is the project-wide standard against
which all abundance arrays are indexed), and each compositional vector is
$\ell^1$-normalised to the unit simplex, $\sum_i\phi_i=1$ --- relative-abundance data
being inherently compositional~\cite{Aitchison1982Compositional, Aitchison1986Compositional, PawlowskyGlahn2015Compositional, Gloor2017Microbiome}. Following
Dieckow et al., a Gaussian-mixture clustering of the W1 compositions splits the
cohort into two \emph{community types}: CT1 ($5$ patients), enriched in early
commensal colonisers (Actinobacteria, Bacilli, Betaproteobacteria), and CT2
($5$ patients), enriched in Fusobacteriia and Bacteroidia. With only
$K\times(T-1)=20$ observations per patient against an interaction matrix of up to
$K(K{+}1)/2=55$ entries, regularisation is not optional.

\begin{table}[htbp]
  \centering\small
  \begin{tabular}{@{}cll@{}}
    \toprule
    \textbf{Index} & \textbf{Guild} & \textbf{SILVA 138.1 class / order}\\
    \midrule
    0 & Actinobacteria      & Actinobacteria, Actinomycetia\\
    1 & Bacilli             & Bacilli\\
    2 & Bacteroidia         & Bacteroidia\\
    3 & Betaproteobacteria  & Betaproteobacteria\\
    4 & Clostridia          & Clostridia, Erysipelotrichia\\
    5 & Coriobacteriia      & Coriobacteriia\\
    6 & Fusobacteriia       & Fusobacteriia, Fusobacteriales\\
    7 & Gammaproteobacteria & Gammaproteobacteria\\
    8 & Negativicutes       & Negativicutes, Selenomonadales, Veillonellales\\
    9 & Other               & All remaining classes\\
    \bottomrule
  \end{tabular}
  \caption[Ten class-level guilds]{The ten class-level guilds (canonical
  ordering). The same guild definitions are used across all datasets in this
  thesis so that abundance arrays and interaction matrices are directly comparable.}
  \label{tab:dieckow_guilds}
\end{table}

%==============================================================================
\section{Three-Layer Metabolic Sign Prior}
\label{sec:dieckow_prior}
%==============================================================================

The sign prior is the thesis's central regulariser. It encodes a single biological
rule --- \emph{if guild $j$ produces a metabolite consumed by guild $i$, then
$A_{ij}>0$ regardless of magnitude} --- and is assembled in three layers of
increasing mechanistic resolution (building on the FBA machinery of
Section~\ref{sec:theory_fba}; Figure~\ref{fig:dieckow_pipeline}).

\paragraph{L1 --- eHOMD / Szafra\'{n}ski experimental flows.} Experimentally
confirmed produces/uses/inhibits relationships between guild-level taxa and oral
metabolites, curated from the expanded Human Oral Microbiome Database
(eHOMD)~\cite{Escapa2018eHOMD} and the Szafra\'{n}ski et al.\ supplement~\cite{Szafranski2025}. Cross-feeding
contributes $+w_{L1}$, competition/inhibition $-w_{L1}$ ($w_{L1}=2.0$ for direct
experimental evidence, $1.5$ literature-supported). L1 alone constrains $34$ guild
pairs.

\paragraph{L2 --- AGORA2 pFBA.} Parsimonious FBA on one representative genome-scale
metabolic model per guild (AGORA2 v2.01~\cite{Heinken2023AGORA2}) under a standardised
oral-fluid medium. pFBA minimises total flux at the optimal growth rate, giving
sparse, sign-robust secretion/uptake patterns; the four strictly anaerobic guilds
have their O$_2$ exchange shut off, and inorganic ions/gases are excluded as
non-syntrophic. Adding L2 ($w_{L2}=1.0$) extends the constrained set only from $34$
to $35$ pairs --- L1 already captures the dominant experimentally-known
cross-feeding.

\paragraph{L3 --- MICOM community FBA.} Single-species pFBA misses the resource
competition and cross-feeding that emerge only when all guilds share one medium.
MICOM~\cite{Diener2020MICOM} solves a community FBA over all $10$ guild models under the
cooperative trade-off
\begin{equation}
  \text{maximise}\;\min_i\!\Bigl(\tfrac{\mu_i}{\mu_i^{\max}}\Bigr)
  \quad\text{s.t.}\quad \textstyle\sum_i\mu_i\geq\tau\sum_i\mu_i^{\max},
  \qquad \tau=0.5,
  \label{eq:dieckow_micom}
\end{equation}
from whose flux solution flux-weighted cross-feeding ($+$) and diffusible-toxin
($-$, e.g.\ H$_2$O$_2$/H$_2$S) signals are extracted. L3 adds $37$ pairs (total
$72$) and the result is robust to $\tau\in\{0.3,0.4,0.5\}$ ($\Delta$RMSE$<0.001$).

\paragraph{Net flow and penalty.} Because the Hamilton matrix is symmetric, the
directed signals are symmetrised, $F_{ij}=\tfrac12(\mathrm{net}[i,j]+\mathrm{net}[j,i])$,
giving prior sign $s_{ij}=\sgn(F_{ij})$ with stiffness $|F_{ij}|$; the asymmetric
gLV uses the directed matrix without symmetrisation. The prior enters the loss as a
half-normal (hinge-squared) penalty that is \emph{zero for correct-sign $A_{ij}$}
and quadratic only for violations,
\begin{equation}
  \calP(\bA;\bF)
  = \sum_{F_{ij}\neq0}\frac{|F_{ij}|}{2\sigma^2}
    \bigl[\max(0,\,-\sgn(F_{ij})\,A_{ij})\bigr]^2,
  \qquad \sigma=0.15,
  \label{eq:dieckow_penalty}
\end{equation}
so it constrains \emph{sign, never magnitude} --- a deliberate choice justified
post hoc in Section~\ref{sec:dieckow_results} (FBA flux units are not ecological
interaction units). The full objective adds a data term and a mild ridge,
$\calL=\mathrm{RMSE}(\hat\bphi,\bphi)+\calP(\bA;\bF)+\lambda\|\bA\|_F^2$
($\lambda=10^{-4}$), minimised by L-BFGS-B with JAX autodiff gradients.

%==============================================================================
\section{Models and Cross-Validation}
\label{sec:dieckow_models}
%==============================================================================

The primary model is the Hamilton replicator (symmetric $\bA$, shared across
patients; patient-specific growth $\bb^{(p)}$) of Section~\ref{sec:theory_replicator},
$\dot\phi_i=\phi_i(f_i(\bphi)-\bar f)$ with $f_i=b_i+\sum_jA_{ij}\phi_j$ and
$A_{ii}\leq0$. For comparison we fit (i) an \emph{asymmetric} gLV (integrated by an explicit Runge--Kutta scheme~\cite{Dormand1980RK45}; no symmetry, no
sign prior unless stated; double the off-diagonal freedom, a ceiling on
expressiveness), (ii) a Hamilton Normal-Species-Pool (NSP) variant optimised by a
three-step implicit-function-theorem scheme, and (iii) MDSINE2~\cite{Bucci2016MDSINE, Gibson2021MDSINE2}, a
fully Bayesian gLV, as a sparsity-induction baseline.

Generalisation is measured by leave-one-patient-out (LOO) cross-validation: for
each held-out patient, $\bA$ and the training $\bb^{(q)}$ are fit on the other
$9$; $\bA$ is frozen; the held-out $\bb^{(p)}$ is fit from W1 only; W2--W3 are
predicted and scored by RMSE and Bray--Curtis (BC) dissimilarity~\cite{Bray1957BrayCurtis, Whittaker1960Diversity}.

%==============================================================================
\section{Results}
\label{sec:dieckow_results}
%==============================================================================

\subsection{Interaction matrix and sign recovery}
\label{sec:dieckow_A}

The fitted matrix (Figure~\ref{fig:dieckow_A}) has all-negative diagonals
(self-limitation) and its dominant facilitations are among early commensal
colonisers --- Actinobacteria--Bacilli ($A_{01}=+1.61$),
Bacilli--Betaproteobacteria ($A_{13}=+1.83$), Actinobacteria--Betaproteobacteria
($A_{03}=+1.67$) --- consistent with documented co-aggregation, while Fusobacteriia
interacts net-negatively with early colonisers (its late-bridging role). The full
L1+L2+L3 prior achieves $70/70$ ($100\%$) sign agreement with the fitted $\bA$ at
$W^*{=}1.0$ --- the lowest weight that achieves full consistency and simultaneously
minimises training RMSE --- and MICOM's $72$ constrained pairs reach $72/72$.

\begin{figure}[htbp]
  \centering
  \includegraphics[width=0.80\textwidth]{dieckow_sign_pipeline.pdf}
  \caption[Three-layer sign-prior construction]{Construction of the metabolic sign
  prior. Oral-fluid medium $\to$ per-guild pFBA $\to$ community MICOM $\to$ signed
  net-flow matrix $F_{ij}$. Constrained pairs by layer: L1 ($34$), L1+L2 ($35$),
  L1+L2+L3 ($72$).}
  \label{fig:dieckow_pipeline}
\end{figure}

\begin{figure}[htbp]
  \centering
  \includegraphics[width=0.72\textwidth]{dieckow_A_matrix.pdf}
  \caption[Fitted in-vivo interaction matrix]{Fitted interaction matrix $\bA$
  (Hamilton, full prior, all $10$ patients). Blue: competition; red: facilitation.
  Diagonals $A_{ii}<0$. Marks: sign agreement (\checkmark) / disagreement ($\times$)
  with the metabolic prior; full-model agreement $70/70$ ($100\%$). The dominant
  red block is the early-coloniser co-aggregation core.}
  \label{fig:dieckow_A}
\end{figure}

The asymmetric gLV ($\alpha{=}0.25$, best model; Table~\ref{tab:dieckow_loo})
recovers the same dominant structure as the Hamilton result, while the directed
representation reveals asymmetric facilitation absent from the symmetric prior
(Figure~\ref{fig:dieckow_glv_A}).

\begin{figure}[htbp]
  \centering
  \includegraphics[width=0.68\textwidth]{dieckow_glv_Amatrix_thesis.pdf}
  \caption[Fitted in-vivo gLV interaction matrix]{Fitted interaction matrix
  $\bA$ (asymmetric gLV, $\alpha{=}0.25$, L1+L2 prior, mean of $10$ LOO folds,
  mean LOO-RMSE $= 0.0490$).
  Blue: competition; red: facilitation. Diagonals $A_{ii}<0$.
  Green borders: sign agreement with metabolic prior (L1+L2);
  red borders: discordant entries. Annotated values shown for $|A_{ij}|>0.05$.
  The early-coloniser co-aggregation block (Actin./Bacil./Bact./$\beta$-Prot.)
  dominates, consistent with the Hamilton result (Figure~\ref{fig:dieckow_A}).}
  \label{fig:dieckow_glv_A}
\end{figure}

The annotated entries reveal three features invisible to the symmetric Hamilton fit.
\emph{Bacillaceae as hub}: the Actin.--Bacil.\ mutual facilitation
($A_{\text{Bacil.,Actin.}}=+3.37$, $A_{\text{Actin.,Bacil.}}=+2.79$) is an order
of magnitude above any other pair; Bacillaceae additionally facilitates
$\beta$-Proteobacteria ($+1.02$), Bacteroidetes ($+1.03$), and Negativicutes
($+0.68$).
\emph{Sign-asymmetric $\gamma$-Proteobacteria}: although Bacillaceae mildly
facilitates $\gamma$-Proteobacteria ($+0.27$), $\gamma$-Proteobacteria strongly
suppresses Bacillaceae in return ($-1.44$, reversed sign) --- a directed antagonism
the symmetric model cannot represent.
\emph{Clostridiaceae broad suppressor}: all ten outgoing interactions are negative
(strongest $A_{\text{Clost.,Bact.}}=-0.66$), a profile that collapses into the
background competition term of the symmetric fit. The strongest diagonal entry is
Actinobacteria self-regulation ($A_{\text{Actin.,Actin.}}=-4.05$).

\begin{figure}[htbp]
  \centering
  \includegraphics[width=0.6\textwidth]{fig_w_sweep_thesis.pdf}
  \caption[AGORA weight sweep: sign agreement and RMSE vs.\ $W$]{Effect of AGORA
  prior weight $W$ on sign agreement (teal, left axis) and training RMSE (orange,
  right axis). $W^*{=}1.0$ is the \emph{first} weight to achieve $100\%$ sign
  agreement ($70/70$ pairs) and simultaneously the one with the lowest training RMSE
  ($0.057$); $W{=}2.0$ also reaches $100\%$ but at higher RMSE ($0.058$). The mean
  LOO-RMSE at $W{=}1.0$ (dotted) confirms that full sign consistency is achieved
  without over-fitting.}
  \label{fig:w_sweep}
\end{figure}

\subsection{Out-of-sample generalisation}
\label{sec:dieckow_loo}

Across all variants (Table~\ref{tab:dieckow_loo}), every mechanistic model
decisively beats the ecological null: the best LOO-BC of $0.147$ is $47.7\%$ below
the persistence null ($0.281$). The sign prior improves the Hamilton model
monotonically (no prior $0.0595\to$ L1+L2 $0.0516\to$ MICOM $0.0502$, the best
Hamilton variant), and the overall best is the asymmetric gLV with $\alpha=0.25$
(L1+L2, RMSE $0.0490\pm0.016$) --- its extra directed freedom slightly outweighs
the symmetric sign constraint for this dataset. MICOM raises sign agreement to a
perfect $72/72$ at a $5$--$6\%$ RMSE cost relative to that best gLV, a
consistency-vs-accuracy trade-off discussed below. The fully Bayesian MDSINE2
collapses to $\bA\approx0$ under only three time points (LOO-RMSE $0.0552$, $13\%$
worse than gLV $\alpha{=}0.25$), confirming that \emph{sign-prior regularisation
outperforms Bayesian sparsity-induction when $T$ is small} --- the operative regime
for clinical microbiome series. The NSP variant ($0.091$) is $1.9\times$ worse,
its latent warm-up cost not repaid on this short horizon.

\begin{table}[htbp]
  \centering\small
  \begin{tabular}{@{}llccc@{}}
    \toprule
    Model & Layer & SA & LOO-RMSE & LOO-BC\\
    \midrule
    Hamilton (no prior)        & ---         & ---            & $0.0595$ & ---\\
    gLV free (asymmetric)      & ---         & ---            & $0.0588$ & $0.154$\\
    Hamilton L1+L2             & L1+L2       & $94\%$         & $0.0516$ & $0.155$\\
    Hamilton +AGORA v1         & L1+L2+L3    & $94\%$         & $0.0562$ & $\mathbf{0.147}$\\
    \textbf{gLV $\alpha{=}0.25$} & \textbf{L1+L2} & $97\%$    & $\mathbf{0.0490}$ & ---\\
    gLV +MICOM ($\alpha{=}0.25$) & L1+L2+L3-MC & $\mathbf{100\%}$ & $0.0516$ & ---\\
    Hamilton +MICOM            & L1+L2+L3-MC & $\mathbf{100\%}$ & $0.0502$ & ---\\
    \midrule
    NSP IFT v2                 & L1+L2       & ---            & $0.0906$ & ---\\
    MDSINE2 (Bayesian gLV)     & Bayesian    & ---            & $0.0552^\dagger$ & ---\\
    \bottomrule
    \multicolumn{5}{@{}l}{\small $^\dagger$8/10 folds; $\bA\approx0$ (insufficient
    Bayesian evidence at $T{=}3$). Null BC: persistence $0.281$.}
  \end{tabular}
  \caption[LOO cross-validation across model variants]{LOO cross-validation. SA:
  sign agreement on constrained pairs. Best values in bold. Sign-prior
  regularisation improves the Hamilton model monotonically; MICOM gives perfect
  sign agreement.}
  \label{tab:dieckow_loo}
\end{table}

\begin{figure}[htbp]
  \centering
  \includegraphics[width=0.80\textwidth]{dieckow_loo_bc.pdf}
  \caption[LOO Bray--Curtis across models]{LOO Bray--Curtis for all variants vs.\
  null models (dashed: persistence $0.281$, grand-mean $0.254$). Every ODE model
  substantially beats both nulls; the full-prior Hamilton attains the lowest BC
  ($0.147$).}
  \label{fig:dieckow_loo}
\end{figure}

\begin{figure}[p]
  \centering
  \includegraphics[width=0.70\textwidth,height=0.50\textheight,keepaspectratio]{dieckow_loo_fits_thesis.pdf}
  \caption[Per-patient LOO predictions: gLV ($\alpha{=}0.25$, L1+L2 prior)]{Per-patient
  leave-one-out predictions for all $10$ Dieckow patients.
  Bars: observed guild fractions at W1 (light), W2 (mid), W3 (dark).
  Crosses ($\times$): LOO-predicted W2 and W3 from the best-model gLV
  ($\alpha{=}0.25$, L1+L2 prior, mean RMSE $= 0.0490$); the held-out patient's
  $\bb^{(p)}$ is fit from W1 only and $\bA$ is frozen from the other 9 patients.
  The model tracks the dominant early-coloniser composition (Actinobacteria,
  Bacilli) across CT1 patients and the Bacteroidia/Fusobacteriia-enriched CT2
  composition, with largest errors at patients with atypical trajectories (K, L).}
  \label{fig:dieckow_loo_fits}
\end{figure}

\subsection{Is mechanistic structure necessary?}
\label{sec:dieckow_benchmark}

Table~\ref{tab:dieckow_loo} compares mechanistic variants among themselves,
but leaves open whether the ecological structure is needed at all. To test this
we benchmarked the gLV model against two baseline families under the identical
LOO one-step-ahead protocol: classical ML regression (Random Forest, Ridge,
LASSO, plus persistence and train-mean references) and modern generative models
(a conditional DDPM and conditional Flow Matching, the ODE-based successor to
diffusion). The mechanistic model wins decisively
(Figure~\ref{fig:dieckow_benchmark}): gLV attains LOO-RMSE $\approx0.05$,
$1.8\times$ below the best ML baseline (Random Forest $0.091$) and $2$--$3\times$
below the generative models (Flow Matching $0.107$, DDPM $0.166$). The gap holds
per patient --- gLV achieves the lowest RMSE on all $10$ held-out patients against
every competitor (one-sided Wilcoxon signed-rank $p=0.001$, the floor at $n=10$).

Two controls localise the cause to the data regime. Flow Matching beats DDPM by
$36\%$ --- the ODE formulation's training-stability advantage over the stochastic
reverse process --- yet both still lose to gLV; and augmenting the ML training set
with diffusion-sampled trajectories \emph{degrades} Ridge and Random Forest,
confirming that with only $18$ real transitions a generative model injects noise
rather than recovering the dynamics. With $T=3$ time points, encoding the
interaction structure is worth more than any amount of model flexibility --- the
same small-$T$ lesson as the MDSINE2 collapse (Table~\ref{tab:dieckow_loo}).

\begin{figure}[htbp]
  \centering
  \includegraphics[width=0.72\textwidth]{dieckow_benchmark.pdf}
  \caption[gLV vs.\ ML and generative baselines]{\textbf{Mechanistic structure
  beats data-driven and generative baselines on LOO prediction.}
  \textbf{(a)}~LOO-CV RMSE (mean $\pm$ SD over the $10$ patients) for the
  one-step-ahead guild-composition task, grouped by family: mechanistic ODE (gLV
  variants), ML regression, and generative models (conditional Flow Matching and
  DDPM). \textbf{(b)}~Paired per-patient comparison of the best model against one
  representative of each competing family; each grey line links the two RMSE
  values for one held-out patient. gLV wins on all $10$ patients against every
  competitor (one-sided Wilcoxon signed-rank $p=0.001$, the minimum attainable at
  $n=10$).}
  \label{fig:dieckow_benchmark}
\end{figure}

\subsection{Sign, not magnitude --- the central finding}
\label{sec:dieckow_signmag}

Two results jointly establish the thesis's organising claim. First, a
\emph{quantitative} Gaussian prior that ties $A_{ij}$ to the FBA flux magnitude
($A_{ij}\sim\mathcal{N}(c\,\Phi_{ij},\sigma^2)$, MacArthur-style) \emph{fails
entirely}, recovering correct signs for only $4$--$11$ of $70$ pairs at any
concentration $c$: FBA flux magnitude is not ecologically predictive. Second, the
\emph{sign}-only prior is both accurate (Table~\ref{tab:dieckow_loo}) and stable:
across the $10$ LOO folds, $31/45$ off-diagonal pairs have perfect sign consistency
(SC$=1.00$), $40/45$ reach SC$\geq0.90$, and only one falls below $0.70$
(Figure~\ref{fig:dieckow_stability}). Prior-pinned pairs keep SC$\geq0.80$ even
when their magnitude is negligible --- the prior fixes direction where the data are
silent. This is why the penalty in Eq.~\eqref{eq:dieckow_penalty} constrains sign
and never magnitude: it injects exactly the part of the metabolic evidence that
transfers (the direction of the interaction) and none of the part that does not
(its flux-unit size).

\begin{figure}[htbp]
  \centering
  \includegraphics[width=0.72\textwidth]{fig_loo_stability.pdf}
  \caption[LOO-CV interaction-sign consistency]{Leave-one-out sign consistency
  of the fitted $\bA$.
  \textbf{(a)}~Fraction of $10$ LOO folds in which each pair $(i,j)$ recovers
  the same sign as the full-data MAP; $31/45$ pairs reach SC$=1.00$, $40/45$
  reach SC$\geq0.90$.
  \textbf{(b)}~Coefficient of variation ($\sigma/|\mu|$) of $A_{ij}$ across
  folds; prior-pinned pairs (dots) remain stable even when magnitudes are small.}
  \label{fig:dieckow_stability}
\end{figure}

A direct test of independence confirms that the cross-feeding directions are
genuine data signal rather than an artefact of the prior.  When the sign
penalty is removed entirely and the symmetric Hamilton model is re-fitted
without guidance, it still recovers $78.6\%$ (11/14 pairs) of the
AGORA-predicted cross-feeding signs --- significantly above a cell-label
permutation null ($\bar p_{\rm null}=37.7\%$,
$p=4\times10^{-4}$, $z=+3.79$, $n_{\rm perm}=10{,}000$).
The asymmetric gLV model shows no such effect ($p=0.96$, $z=-1.47$):
the Hamilton symmetry constraint mediates the metabolic--ecological
correspondence.  Only cooperative (cross-feeding) directions are validated;
competitive interactions carry no detectable signal across patients,
consistent with competition being shaped by local context rather than
conserved metabolic architecture.

\subsection{Patient structure and community type}
\label{sec:dieckow_struct}

The shared $\bA$ captures patient-independent guild ecology; patient-level
variation is absorbed by the growth vectors $\bb^{(p)}$. A PCA of the fitted
$\hat\bb^{(p)}$ (Figure~\ref{fig:dieckow_bpca}) separates CT1 from CT2 along PC1
\emph{without} community-type supervision (CT1: higher Actinobacteria/Bacilli
intrinsic rates; CT2: higher Fusobacteriia), and the LOO predictions preserve the
CT1/CT2 discriminant (LDA accuracy $90\%$, only the atypical Patient~A
misclassified) while spanning the observed compositional manifold rather than
collapsing to the mean. Off-diagonal interactions are on average twice as strong as
intrinsic rates ($\overline{|A_{\mathrm{off}}|}=0.35$ vs.\ $\overline{|b|}=0.18$):
community ecology dominates individual guild growth here.


\subsection{Network centrality and keystone structure}
\label{sec:dieckow_network}

Reading the fitted $\bA$ as a directed weighted graph~\cite{Freilich2020Inference, May1972Stability}
exposes which guilds organise the community (Figure~\ref{fig:dieckow_centrality}). \emph{Actinobacteria} is the
dominant hub on every measure: highest outgoing influence ($5.6$), the largest
positive net effect (out minus in, $+2.0$), and top centrality rank --- a net
\emph{source} that drives the early-coloniser core. \emph{Bacilli} is the most
\emph{vulnerable} guild (incoming strength $8.2$): abundant but strongly
acted-upon. The most ecologically interesting guild is \emph{Fusobacteriia}, which
carries a large \emph{keystone gap} (centrality rank well above its abundance rank,
gap $+2$): it is influential out of proportion to how much of it there is --- the
quantitative signature of a low-abundance bridging organism, the in-vivo guild-level
echo of the \textit{F.~nucleatum} bridging role seen in vitro
(Section~\ref{sec:heine_rewiring}). These rankings are stable across the LOO folds,
so the keystone ordering is a property of the inferred ecology, not of any single
patient.

\begin{figure}[htbp]
  \centering
  \includegraphics[width=0.70\textwidth]{dieckow_centrality.pdf}
  \caption[Guild network centrality]{Guild interaction-network centrality from the
  fitted $\bA$. Actinobacteria is the top influencer / net source; Bacilli is the
  most vulnerable (highest incoming strength); Fusobacteriia carries the largest
  keystone gap (central despite low abundance --- a bridging guild).}
  \label{fig:dieckow_centrality}
\end{figure}

\begin{figure}[htbp]
  \centering
  \includegraphics[width=0.70\textwidth]{fig_network_backbone_thesis.pdf}
  \caption[Cross-cohort concordant interaction backbone]{Concordant interaction
  backbone across Dieckow ($N=10$) and Botelho ($N=15$) cohorts.
  Strong pairs ($|\hat A_{ij}|>0.1$ in both fits) are shown; $8/9$ agree in sign
  ($p\approx0.02$, permutation). The Actinobacteria hub is the most concordant
  node (5/5 edges consistent), confirming its role as a structural hub independent
  of cohort.}
  \label{fig:dieckow_backbone}
\end{figure}

\begin{figure}[htbp]
  \centering
  \includegraphics[width=0.70\textwidth]{fig_network_rewiring_thesis.pdf}
  \caption[CS$\to$DH interaction network rewiring]{Network rewiring from
  Commensal Static (CS) to Dysbiotic HOBIC (DH) in the Heine 5-species system.
  Edges that gain strength ($|\Delta A_{ij}|>0.2$) are highlighted; the
  \textit{Pg}--\textit{Fn} positive edge and the competitive suppression of
  \textit{An}/\textit{So} are the dominant rewired pairs.}
  \label{fig:dieckow_rewiring}
\end{figure}

%==============================================================================
\section{External Validation: Peri-Implant Dysbiosis}
\label{sec:dieckow_joshi}
%==============================================================================

To test transfer beyond the calibration cohort, the Dieckow-derived dysbiosis
ordering is applied to $127$ independent cross-sectional peri-implant samples
(Joshi et al.~\cite{joshi2025}: Health/Mucositis/Peri-implantitis). A guild-level
red/green dysbiosis ratio,
$\log(\phi_{\mathrm{Fuso}}+\phi_{\mathrm{Bact}}+\phi_{\mathrm{Clos}})-
\log(\phi_{\mathrm{Bac}}+\phi_{\mathrm{Act}}+\phi_{\mathrm{Neg}})$,
cleanly separates the diagnostic groups (Table~\ref{tab:dieckow_joshi};
Figure~\ref{fig:dieckow_joshi}): Health and Mucositis sit together (mean $-2.09$
vs.\ $-1.74$) while Peri-implantitis shifts sharply positive ($+0.53$).
Kruskal--Wallis $H=34.7$, $p<10^{-4}$; severity correlation Spearman $\rho=0.46$
($p=5\times10^{-8}$, $N=127$), rising to $\rho=0.59$ for the binary
Health-vs-PI contrast. The model thus orders dysbiosis correctly on data it never
saw; the moderate three-way $\rho$ reflects the near-identity of Health and
Mucositis in this five-genus space, which would need additional pathogens to
resolve.

\begin{table}[htbp]
  \centering\small
  \begin{tabular}{@{}lcccc@{}}
    \toprule
    Diagnosis & $n$ & Mean R/G & Median R/G & vs.\ Health\\
    \midrule
    Health           & $56$ & $-2.09$ & $-1.88$ & ---\\
    Mucositis        & $39$ & $-1.74$ & $-1.65$ & $p=0.47$\\
    Peri-implantitis & $32$ & $+0.53$ & $+0.76$ & $p<10^{-4}$\\
    \bottomrule
    \multicolumn{5}{@{}l}{\small Kruskal--Wallis $H=34.7$, $p<10^{-4}$; Spearman
    $\rho=0.46$, $p=5\times10^{-8}$.}
  \end{tabular}
  \caption[Peri-implant dysbiosis ordering]{Guild-level R/G dysbiosis ratio by
  diagnosis (Joshi et al., $127$ samples). The Dieckow-derived ordering separates
  peri-implantitis from health/mucositis on independent cross-sectional data.}
  \label{tab:dieckow_joshi}
\end{table}

\begin{figure}[htbp]
  \centering
  \includegraphics[width=0.68\textwidth,height=0.36\textheight,keepaspectratio]{dieckow_joshi.pdf}
  \caption[Peri-implant dysbiosis attractor analysis]{R/G dysbiosis ratio across
  Health, Mucositis and Peri-implantitis (Joshi et al., $N=127$). Peri-implantitis
  separates with a large positive shift; severity correlation $\rho=0.46$
  ($p=5\times10^{-8}$).}
  \label{fig:dieckow_joshi}
\end{figure}

\begin{figure}[htbp]
  \centering
  \includegraphics[width=0.72\textwidth]{fig_prior_asymmetry.pdf}
  \caption[Prior sign agreement by community state]{Sign agreement between the
  AGORA-predicted cross-feeding signs and the TMCMC posterior, with and without
  the metabolic prior, across the four Heine attractors.  Health states (CS, CH)
  gain $+24$ and $+56$ percentage points from the prior; dysbiotic states (DS, DH)
  lose $-44$ and $-10$ pp, falling below chance.  The metabolic prior is
  \emph{informative for health-associated ecology and actively harmful for
  dysbiotic ecology}, where pathogenic interactions override the commensal
  cross-feeding template.}
  \label{fig:prior_asymmetry}
\end{figure}

%==============================================================================
\section{Discussion}
\label{sec:dieckow_discussion}
%==============================================================================

The in-vivo analysis confirms and sharpens the thesis's organising idea. Metabolic
flux-balance evidence is predictive \emph{as a sign constraint but not as a
magnitude} (the quantitative MacArthur prior fails outright), and a sign-only prior
is what generalises: it improves out-of-sample prediction monotonically, stabilises
the inferred network across folds, and lets community FBA (MICOM) reach perfect
sign agreement. That MICOM's full consistency comes at a small RMSE cost relative
to the best gLV reflects the honest tension between biological self-consistency and
raw predictive accuracy at this cohort size, not a failure of either. The framework
also out-predicts a fully Bayesian gLV (MDSINE2), which collapses to no interactions
under three time points --- the practically important lesson that mechanistic sign
priors substitute for the data a Bayesian model would otherwise need.

The principal limitations are statistical and structural. With $N=10$ patients the
AGORA/MICOM improvement over L1+L2 is real but only marginally significant
($p\approx0.07$--$0.08$), so the perfect-sign-agreement result should be read as a
consistency check rather than a powered effect. The Hamilton symmetry constraint
$A_{ij}=A_{ji}$ excludes genuinely asymmetric ecological interactions --- precisely
where the asymmetric gLV gains its edge --- and the prior is curated for a
caries-recovery context whose transfer to other disease states is untested beyond
the peri-implant ordering shown above. These motivate the spatial extension of
Chapter~\ref{ch:integration}, where the same prior-regularised interactions are
carried into a depth-resolved reaction--diffusion model and confronted with FISH
imaging.
A further limitation is that the magnitude of AGORA2-derived cross-feeding
fluxes cannot be directly calibrated to the gLV interaction scale without
additional kinetic parameters (Monod half-saturation constants); recovering
absolute interaction strengths from GEM stoichiometry remains an open problem
in the field~\cite{Rath2024GEMReliability,Mustri2025GLVABreakdown}.

\paragraph{Statistical power and the sign enrichment as primary evidence.}
The marginal significance of the MICOM--gLV RMSE comparison ($p\approx0.07$,
$N=10$) is expected: for a medium effect size ($d=0.5$) and $N=10$ a
two-sample test has power $\approx0.30$, so the study is structurally
underpowered for this particular comparison. The primary quantitative
evidence is therefore the permutation-based cross-feeding sign enrichment
reported above, which aggregates over all $45\times10$ fold-pair combinations
and is independent of cohort size in this different sense. A complementary
robustness check --- neutral-initialised no-prior LOO starting from
$\mathbf{A}=\mathbf{0}$, which eliminates any warm-start bias --- yields
$\mathrm{SA}=71.4\%$ ($p=0.027$, permutation test, $n_\pi=10^4$), confirming
that the cross-feeding enrichment is data-driven and not an artefact of
warm-starting from the prior-fitted $\mathbf{A}$.
